\documentclass[worksheet]{workbook}

\solutionfalse
% \solutiontrue  % uncomment to reveal solutions

\assignmenttitle{Practice Worksheet 3: Concept Maps}
% \assignmenttitle{Solutions for Practice Worksheet 2}
\assignmentcourse{\textsc{Math 163}}
\assignmentsemester{Spring 2026}
\assignmentduedate{Math 163}
\instructorname{Prof.\ Chowdhury}

\begin{document}

\maketitle

The next page displays fifteen major concepts from chapter 2 as labeled nodes, connected by numbered arrows.  Each arrow represents a mathematical relationship: a \textbf{solid arrow} means the source concept implies or establishes the target; a \textbf{dashed arrow} marks a relationship that \alert{fails}; a \textbf{double-headed arrow} means the two concepts are equivalent.  Your task: \textit{for each numbered arrow, write a one-sentence description of the relationship it represents.}  Name the relevant theorem or give a precise mathematical reason.

\bigskip

\noindent{\large\bfseries\color{maroon} Concept Inventory}

\medskip
\small

\begin{enumerate}[label=\textbf{\Alph*.}, leftmargin=2.5em, itemsep=1pt]

\item \textbf{Infinite Series $\sum a_n$.}  An expression $a_1 + a_2 + a_3 + \cdots$ whose convergence is defined via its partial sums.

\item \textbf{Partial Sums $\{s_n\}$.}  $s_n = a_1 + a_2 + \cdots + a_n$. The series converges iff $\{s_n\}$ converges.

\item \textbf{Convergent Series.}  The series $\sum a_n$ converges when $\{s_n\}$ has a finite limit.

\item \textbf{Divergent Series.}  The series $\sum a_n$ diverges when $\{s_n\}$ fails to converge.

\item \textbf{Divergence Test.}  If $a_n \not\to 0$, then $\sum a_n$ diverges.  The converse is false ($\sum 1/n$).

\item \textbf{Geometric Series $\sum ar^n$.}  Converges iff $|r| < 1$, sum $= a/(1-r)$.  Closed-form partial sums: $s_n = a(1-r^n)/(1-r)$.

\item \textbf{Telescoping Series.}  $\sum (b_n - b_{n+1})$ converges iff $\{b_n\}$ converges, with sum $= b_1 - \lim b_n$.  Partial sums collapse.

\item \textbf{$p$-Series $\sum 1/n^p$.}  Converges iff $p > 1$.  The boundary $p = 1$ is the harmonic series.  Proved via the integral test.

\item \textbf{Nonneg Series $+$ MCT.}  If $a_n \geq 0$, the partial sums are increasing.  By the MCT, $\sum a_n$ converges iff $\{s_n\}$ is bounded.  Foundation for all nonneg tests.

\item \textbf{Comparison Tests.}  Direct: $0 \leq a_n \leq cb_n$ eventually and $\sum b_n$ converges $\implies$ $\sum a_n$ converges.  Limit: $\lim a_n/b_n \in (0, \infty) \implies$ same behavior.

\item \textbf{Integral Test.}  If $f$ is positive, continuous, decreasing, and $f(n) = a_n$, then $\sum a_n$ and $\int_1^\infty f$ converge or diverge together.

\item \textbf{Dirichlet's Test.}  If $\{a_n\}$ decreases to $0$ and the partial sums $B_n = \sum_{k=1}^n b_k$ are bounded, then $\sum a_n b_n$ converges.  Alternating series test: special case $b_n = (-1)^n$.  Proof: summation by parts.

\item \textbf{Absolute Convergence.}  $\sum a_n$ converges absolutely if $\sum |a_n|$ converges.  Absolute convergence implies convergence (via Cauchy criterion).

\item \textbf{Conditional Convergence.}  $\sum a_n$ converges conditionally if it converges but $\sum |a_n|$ diverges.  Depends on cancellation.

\item \textbf{Ratio and Root Tests.}  $L = \lim |a_{n+1}/a_n|$ or $L = \limsup |a_n|^{1/n}$: if $L < 1$, absolutely convergent; if $L > 1$, divergent; if $L = 1$, inconclusive.  Both compare with geometric series.

\end{enumerate}

\normalsize

\medskip

\noindent\textbf{Other tools and benchmarks:} Cauchy criterion for series; harmonic series $\sum 1/n$ (diverges despite $a_n \to 0$); alternating harmonic $\sum (-1)^{n+1}/n$ (conditionally convergent, sum $= \ln 2$); summation by parts (discrete integration by parts, key tool for Dirichlet's test); $\sum \sin(n\theta)/n$ (non-alternating Dirichlet example).


% ============================================================
% PAGE 2: CONCEPT MAP DIAGRAM
% ============================================================
\newpage
\begin{center}
{\large\bfseries\color{maroon} Concept Map}
\end{center}

\begin{center}
\begin{tikzpicture}[>=Latex,
  concept/.style={rectangle, draw=darkgray, rounded corners=2pt, minimum width=2.2cm, minimum height=0.6cm, align=center, font=\footnotesize, inner sep=3pt, fill=white, line width=0.5pt},
  highlight/.style={concept, fill=lightblue!15},
  edgenum/.style={circle, fill=white, draw=darkgray, inner sep=0pt, minimum size=11pt, font=\scriptsize\bfseries, line width=0.4pt},
  arr/.style={->, thick, darkgray},
  fail/.style={->, thick, dashed, red!70!black},
  nlabel/.style={font=\tiny\bfseries, text=darkgray, inner sep=1pt}
]

% ===== ROW 0: Definition =====
\node[concept] (A) at (1.5, 0)
  {Infinite Series\\[-1pt]$\sum a_n$};

% ===== ROW 1: Partial sums, Divergence test =====
\node[highlight] (B) at (-1.5, -2)
  {Partial Sums\\[-1pt]$\{s_n\}$};
\node[concept] (E) at (6, -2)
  {Divergence\\[-1pt]Test};

% ===== ROW 2: Benchmarks with closed-form sums =====
\node[concept] (F) at (-1, -3.8)
  {Geometric\\[-1pt]Series};
\node[concept] (G) at (3, -3.8)
  {Telescoping\\[-1pt]Series};

% ===== ROW 3: Convergent / Divergent =====
\node[concept, fill=lightgreen!12] (C) at (-4, -5.8)
  {Convergent\\[-1pt]Series};
\node[concept] (D) at (7, -5.8)
  {Divergent\\[-1pt]Series};

% ===== ROW 4: Main test families =====
\node[concept] (I) at (-6.5, -8.5)
  {Nonneg Series\\[-1pt]$+$ MCT};
\node[concept] (L) at (-2, -8.5)
  {Dirichlet's\\[-1pt]Test};
\node[concept] (M) at (2.5, -8.5)
  {Absolute\\[-1pt]Convergence};
\node[concept] (O) at (7, -8.5)
  {Ratio and\\[-1pt]Root Tests};

% ===== ROW 5: Bottom tools =====
\node[concept] (J) at (-7.5, -11)
  {Comparison\\[-1pt]Tests};
\node[concept] (K) at (-4, -11)
  {Integral\\[-1pt]Test};
\node[concept] (H) at (-0.5, -11)
  {$p$-Series\\[-1pt]$\sum 1/n^p$};
\node[concept] (N) at (4, -11)
  {Conditional\\[-1pt]Convergence};

% ===== NODE LETTER LABELS =====
% Inventory letters: A=A, B=B, C=C, D=D, E=E, F=F, G(telescoping)=G, 
% H(p-series)=H, I(nonneg)=I, J(comparison)=J, K(integral)=K,
% L(Dirichlet)=L, M(absolute)=M, N(conditional)=N, O(ratio/root)=O
\node[nlabel, anchor=south east] at (A.north west) {A};
\node[nlabel, anchor=south east] at (B.north west) {B};
\node[nlabel, anchor=south east] at (C.north west) {C};
\node[nlabel, anchor=south east] at (D.north west) {D};
\node[nlabel, anchor=south east] at (E.north west) {E};
\node[nlabel, anchor=south east] at (F.north west) {F};
\node[nlabel, anchor=south east] at (G.north west) {G};
\node[nlabel, anchor=south east] at (H.north west) {H};
\node[nlabel, anchor=south east] at (I.north west) {I};
\node[nlabel, anchor=south east] at (J.north west) {J};
\node[nlabel, anchor=south east] at (K.north west) {K};
\node[nlabel, anchor=south east] at (L.north west) {L};
\node[nlabel, anchor=south east] at (M.north west) {M};
\node[nlabel, anchor=south east] at (N.north west) {N};
\node[nlabel, anchor=south east] at (O.north west) {O};


% ================================================================
%  EDGES — 16 numbered, unlabeled
% ================================================================

% 1: A → B  (convergence defined via partial sums)
\draw[arr] (A) -- (B)
  node[edgenum, pos=0.45, above left=-2pt] {1};

% 2: B → C  (partial sums converge → series converges)
\draw[arr] (B.south west) -- (C.north)
  node[edgenum, pos=0.5, left=2pt] {2};

% 3: B → D  (partial sums diverge → series diverges)
\draw[arr] (B.east) to[out=-10, in=150]
  node[edgenum, pos=0.3, above=3pt] {3} (D.north west);

% 4: E → D  (divergence test: a_n ≠→ 0 forces divergence)
\draw[arr] (E) -- (D)
  node[edgenum, midway, right=2pt] {4};

% 5: E ⇸ C  (a_n → 0 necessary but NOT sufficient)
%   Route: from E going left-and-down, ending with × above the C/D row
\draw[fail] (E.south west) -- (C.north east)
  node[edgenum, pos=0.2, above left=-1pt] {5}
  node[at end, font=\normalsize\bfseries, red!70!black] {$\boldsymbol{\times}$};

% 6: F → C  (geometric: converges iff |r| < 1)
\draw[arr] (F.south west) -- (C.north)
  node[edgenum, pos=0.5, above left=-2pt] {6};

% 7: G → C  (telescoping: partial sums collapse)
\draw[arr] (G.south west) to[out=210, in=30]
  node[edgenum, pos=0.45, above=2pt] {7} (C.east);

% 8: I → C  (MCT: bounded nonneg partial sums → convergence)
\draw[arr] (I.north) to[out=80, in=220]
  node[edgenum, pos=0.5, left=2pt] {8} (C.south west);

% 9: I → J  (comparison tests flow from MCT)
\draw[arr] (I.south) -- (J.north)
  node[edgenum, pos=0.45, left=2pt] {9};

% 10: I → K  (integral test flows from MCT)
\draw[arr] (I.south east) -- (K.north)
  node[edgenum, pos=0.45, right=2pt] {10};

% 11: K → H  (integral test determines p-series threshold)
\draw[arr] (K) -- (H)
  node[edgenum, midway, above] {11};

% 12: L → C  (Dirichlet gives convergence; alternating as corollary)
\draw[arr] (L.north) -- (C.south)
  node[edgenum, pos=0.45, right=2pt] {12};

% 13: M → C  (absolute convergence implies convergence)
\draw[arr] (M.north west) to[out=135, in=325]
  node[edgenum, pos=0.5, right=2pt] {13} (C.south east);

% 14: C ⇸ M  (convergence does NOT imply absolute)
\draw[fail] (C.south) -- (M.south west)
  node[edgenum, pos=0.5, below left=-1pt] {14}
  node[at end, font=\normalsize\bfseries, red!70!black] {$\boldsymbol{\times}$};

% 15: O → M  (ratio/root establish absolute convergence)
\draw[arr] (O.west) -- (M.east)
  node[edgenum, midway, above] {15};

% 16: L → N  (Dirichlet can produce conditional convergence)
\draw[arr] (L.south east) to[out=-30, in=150]
  node[edgenum, pos=0.45, above=2pt] {16} (N.north west);

\end{tikzpicture}
\end{center}

\medskip

\noindent{\bfseries Edge Labels.} For each numbered arrow, describe the relationship.  Name the theorem, state the key hypothesis, or explain what goes wrong.

\smallskip
\footnotesize
\begin{multicols}{2}
\begin{enumerate}[label=\textbf{\arabic*.}, itemsep=0.8em, leftmargin=2em]
\item \dotfill
\item \dotfill
\item \dotfill
\item \dotfill
\item \dotfill
\item \dotfill
\item \dotfill
\item \dotfill
\columnbreak
\item \dotfill
\item \dotfill
\item \dotfill
\item \dotfill
\item \dotfill
\item \dotfill
\item \dotfill
\item \dotfill
\end{enumerate}
\end{multicols}

\normalsize

\noindent\textbf{Reflection.} Which single arrow represents the most important insight of this chapter?  Explain in one sentence.

\vfill













\newpage


% ============================================================
%  CONCEPT MAP ACTIVITY — append at end of Chapter 3
%  (before any \end{document} or next chapter)
% ============================================================

\newpage
\begin{center}
{\Large\bfseries\color{maroon} Sequences and Series of Functions}
\end{center}

\medskip

\noindent The next page displays fourteen major concepts from chapter 3 as labeled nodes, connected by numbered arrows.  Each arrow represents a mathematical relationship: a \textbf{solid arrow} means the source concept implies or establishes the target; a \textbf{dashed red arrow} marks a relationship that \alert{fails}; a \textbf{double-headed arrow} means equivalence.  One arrow is drawn in \textbf{\textcolor{orange!80!black}{dotted orange}}: it marks a theorem whose hypotheses differ from those of its neighbors.  Your task: \textit{for each numbered arrow, write a one-sentence description of the relationship it represents.}  Name the relevant theorem, state the key hypotheses, or explain what goes wrong.

\bigskip

\noindent{\large\bfseries\color{maroon} Concept Inventory}

\medskip
\small

\begin{enumerate}[label=\textbf{\Alph*.}, leftmargin=2.5em, itemsep=1pt]

\item \textbf{Pointwise Convergence of $\{f_n\}$.}  For each $x$ and each $\epsilon > 0$, there exists $N$ (depending on \emph{both} $x$ and $\epsilon$) such that $n \geq N \implies |f_n(x) - f(x)| < \epsilon$.

\item \textbf{Uniform Convergence of $\{f_n\}$.}  For each $\epsilon > 0$, there exists $N$ (depending only on $\epsilon$) such that $n \geq N \implies |f_n(x) - f(x)| < \epsilon$ for \emph{all} $x$ simultaneously.

\item \textbf{Cauchy Criterion for Uniform Convergence.}  $\{f_n\}$ converges uniformly iff for each $\epsilon > 0$ there exists $N$ such that $|f_m(x) - f_n(x)| < \epsilon$ for all $m, n \geq N$ and all $x$.

\item \textbf{Continuity of the Limit.}  Each $f_n$ continuous $+$ $f_n \unifto f$ $\implies$ $f$ continuous.

\item \textbf{Integrability of the Limit.}  Each $f_n$ integrable $+$ $f_n \unifto f$ $\implies$ $\lim \int f_n = \int \lim f_n$.

\item \textbf{Differentiability of the Limit \emph{(stronger hypotheses)}.}  $f_n \to f$ \textbf{pointwise} $+$ $f_n' \unifto g$ \textbf{uniformly} $\implies$ $f' = g$.  Note: it is the \emph{derivatives} that must converge uniformly, not $f_n$ itself.  Uniform convergence of $f_n$ alone does not suffice.

\item \textbf{Dini's Theorem.}  $f_n \to f$ pointwise on $[a,b]$, convergence monotone, each $f_n$ and $f$ continuous $\implies$ $f_n \unifto f$.

\item \textbf{Series of Functions $\sum f_n$.}  Convergence means the partial sums $S_N(x) = \sum_{k=1}^N f_k(x)$ converge.  A series of functions \emph{is} a sequence of functions (the partial sums).

\item \textbf{Weierstrass $M$-test.}  If $|f_n(x)| \leq M_n$ for all $x$ and $\sum M_n < \infty$, then $\sum f_n$ converges uniformly and absolutely.

\item \textbf{Power Series $\sum a_n(x-c)^n$.}  A series of functions whose $n$th term is a monomial of degree $n$ centered at $c$.

\item \textbf{Radius of Convergence $R$.}  $\sum a_n(x-c)^n$ converges absolutely for $|x-c| < R$, diverges for $|x-c| > R$.  Endpoints checked separately.

\item \textbf{Term-by-term Differentiation and Integration.}  Inside $(c-R, c+R)$, a power series can be differentiated and integrated term by term.  The derived series has the same radius $R$.

\item \textbf{Taylor Series.}  $\sum_{n=0}^\infty \frac{f^{(n)}(a)}{n!}(x-a)^n$.  The Taylor coefficients are forced by the differentiability of the power series: $a_n = f^{(n)}(c)/n!$.

\item \textbf{Convergence of Taylor Series to $f$.}  The Taylor series converges to $f(x)$ iff the remainder $R_N(x) = f(x) - P_N(x) \to 0$.

\end{enumerate}

\normalsize

\medskip

\noindent\textbf{Key counterexamples:} $x^n$ on $[0,1]$ (continuous $\to$ discontinuous); $\sqrt{x^2+1/n^2}\unifto |x|$ (differentiable $\to$ non-differentiable); $(1/n)\sin(n^2 x)$ (derivatives diverge); $x/(1+n^2 x^2)\unifto 0$ ($ \ddx\lim \neq \lim \ddx$); $e^{-1/x^2}$ (smooth but not analytic).


% ============================================================
% PAGE 2: CONCEPT MAP DIAGRAM
% ============================================================
\newpage
\begin{center}
{\large\bfseries\color{maroon} Concept Map}
\end{center}

\begin{center}
\begin{tikzpicture}[>=Latex,
  concept/.style={rectangle, draw=darkgray, rounded corners=2pt, minimum width=2.3cm, minimum height=0.65cm, align=center, font=\footnotesize, inner sep=3pt, fill=white, line width=0.5pt},
  highlight/.style={concept, fill=lightblue!15},
  edgenum/.style={circle, fill=white, draw=darkgray, inner sep=0pt, minimum size=12pt, font=\scriptsize\bfseries, line width=0.4pt},
  arr/.style={->, thick, darkgray},
  darr/.style={<->, thick, darkgray},
  fail/.style={->, thick, dashed, red!70!black},
  caution/.style={->, thick, orange!80!black, densely dotted},
  nlabel/.style={font=\tiny\bfseries, text=darkgray, inner sep=1pt}
]

% ===== ROW 0: Types of convergence =====
\node[concept] (A) at (-6, 0)
  {Pointwise\\[-1pt]Convergence};
\node[highlight] (B) at (1.5, 0)
  {Uniform\\[-1pt]Convergence};
\node[concept] (C) at (7.5, 0.8)
  {Cauchy\\[-1pt]Criterion};

% ===== ROW 1: Interchange theorems =====
\node[concept] (D) at (-0.5, -3.3)
  {Continuity\\[-1pt]of Limit};
\node[concept] (E) at (3.5, -3.3)
  {Integrability\\[-1pt]of Limit};
\node[concept, draw=orange!80!black] (F) at (7.5, -3.3)
  {Differentiability\\[-1pt]of Limit};

% ===== ROW 2: Tools + bridge to series =====
\node[concept] (G) at (-5.5, -6.3)
  {Dini's\\[-1pt]Theorem};
\node[concept, fill=lightgreen!12] (H) at (-1.5, -6.3)
  {Series of Functions\\[-1pt]$\sum f_n$ via $\{S_N\}$};
\node[concept] (I) at (4, -6.3)
  {Weierstrass\\[-1pt]$M$-test};

% ===== ROW 3: Power series =====
\node[concept] (J) at (0, -9)
  {Power Series\\[-1pt]$\sum a_n(x-c)^n$};
\node[concept] (K) at (6, -9)
  {Radius of\\[-1pt]Convergence $R$};

% ===== ROW 4: Taylor =====
\node[concept] (L) at (-2, -11.3)
  {Term-by-term\\[-1pt]Diff.\ and Integ.};
\node[concept] (M) at (3, -11.3)
  {Taylor\\[-1pt]Series};
\node[concept] (N) at (7.5, -11.3)
  {Taylor series\\[-1pt]converges to $f$};


% ===== NODE LETTER LABELS =====
\foreach \nd/\lbl in {A/A, B/B, C/C, D/D, E/E, F/F, G/G, H/H, I/I, J/J, K/K, L/L, M/M, N/N} {
  \node[nlabel, anchor=south east] at (\nd.north west) {\lbl};
}


% ================================================================
%  EDGES — 15 numbered, unlabeled
% ================================================================

% 1: B → A  (uniform implies pointwise)
\draw[arr] (B) -- (A)
  node[edgenum, midway, above] {1};

% 2: A ⇸ D/E/F  (pointwise does NOT preserve properties)
\draw[fail] (A.south) -- (D.west)
  node[edgenum, pos=0.3, left=2pt] {2}
  node[at end, font=\normalsize\bfseries, red!70!black] {$\boldsymbol{\times}$};

% 3: B → D  (Continuous Uniform Limit Theorem)
\draw[arr] (B.south west) -- (D.north)
  node[edgenum, pos=0.5, above left=-2pt] {3};

% 4: B → E  (Integrable Uniform Limit Theorem)
\draw[arr] (B.south) -- (E.north)
  node[edgenum, pos=0.5, right=2pt] {4};

% 5: B → F  ⚠ (Differentiable ULT — needs f_n' →unif g, NOT f_n →unif f)
%   Drawn in orange dotted to signal different hypothesis structure
\draw[caution] (B.south east) -- (F.north)
  node[edgenum, draw=orange!80!black, pos=0.4, right=2pt] {5}
  node[pos=0.6, right=5pt, font=\scriptsize\bfseries, orange!80!black] {$\boldsymbol{!}$};

% 6: C ↔ B  (equivalent characterization)
\draw[darr] (B.east) -- (C.west)
  node[edgenum, midway, above] {6};

% 7: G → B  (Dini: ptwise + monotone + cont + compact ⟹ uniform)
\draw[arr] (G.north) to[out=45, in=215]
  node[edgenum, pos=0.75, left=2pt] {7} (B.south west);

% 8: H → A  (series conv = convergence of partial sums = sequence conv)
\draw[arr] (H.north west) -- (A.south)
  node[edgenum, pos=0.5, left=2pt] {8};

% 9: I → B  (M-test establishes uniform convergence for series)
\draw[arr] (I.west) to[out=180, in=270]
  node[edgenum, pos=0.55, right=2pt] {9} (B.south);

% 10: J → K  (trichotomy theorem determines R)
\draw[arr] (J.east) -- (K.west)
  node[edgenum, midway, above] {10};

% 11: J → B  (power series converge uniformly on compact sub-intervals)
\draw[arr] (J.north) -- (B.south)
  node[edgenum, pos=0.45, left=3pt] {11} ;

% 12: J → L  (term-by-term operations via uniform conv + interchange)
\draw[arr] (J.south) -- (L.north)
  node[edgenum, pos=0.45, above left=-2pt] {12};

% 13: L → M  (differentiability gives Taylor coefficients: a_n = f^(n)(c)/n!)
\draw[arr] (L) -- (M)
  node[edgenum, midway, above] {13};

% 14: M → N  (Lagrange remainder criterion → convergence to f)
\draw[arr] (M) -- (N)
  node[edgenum, midway, above] {14};

% 15: H → D  (series corollaries: M-test + interchange → 
%              continuity/integrability/differentiability of sums)
\draw[arr] (H.north) -- (D.south)
  node[edgenum, pos=0.5, above=2pt] {15};

\end{tikzpicture}
\end{center}

\medskip

\noindent{\bfseries Edge Labels.} For each numbered arrow, describe the relationship.  Name the theorem, state the key hypothesis, or explain what goes wrong.  Pay special attention to arrow~5.

\smallskip
\footnotesize
\begin{multicols}{2}
\begin{enumerate}[label=\textbf{\arabic*.}, itemsep=0.85em, leftmargin=2em]
\item \dotfill
\item \dotfill
\item \dotfill
\item \dotfill
\item \dotfill
\item \dotfill
\item \dotfill
\item \dotfill
\columnbreak
\item \dotfill
\item \dotfill
\item \dotfill
\item \dotfill
\item \dotfill
\item \dotfill
\item \dotfill
\end{enumerate}
\end{multicols}

\normalsize

\noindent\textbf{Reflection.} Which single arrow represents the most important insight of this chapter?  Explain in one sentence.

\vfill

\end{document}